\chapter{Results and Evaluation}

\section{Experiment Setup}

The experiment evaluates four extraction strategies against a
hand-written ground truth. The ground truth consists of 90 Markdown
files written manually by the thesis author, 30 per thesis document.
Each file was produced by reading one page of the original PDF and
writing the Markdown that a correct conversion should produce.
Headings were assigned the correct depth from the document's numbering
hierarchy, code blocks were wrapped in fenced code blocks with the
appropriate language tag, tables were reconstructed in pipe-table
syntax, and formulas were written in LaTeX notation.
Writing these files by hand was necessary to ensure an honest ground
truth: an automatically generated reference would contain the same
errors as the system under evaluation, making the comparison circular.

The three documents in the corpus are a German computer science
bachelor thesis on neural networks (referred to here as the neural
networks thesis), a German bachelor thesis on sentiment analysis
(the sentiment analysis thesis), and an English-language bachelor
thesis describing a human-in-the-loop system for research paper
generation (the HITL thesis). Together they cover different writing
styles, content densities, and proportions of code, formula, and
table content, which makes the per-page results more representative
than a single document would allow.

All four strategies (text-only, image-only, hybrid, and adaptive) were
applied to every page of every document. This produced
$4 \times 3 \times 30 = 360$ evaluated conversions with zero errors.
For each conversion, the pipeline output was compared to the
corresponding reference file using the seven metrics described in
Section~4.3: text similarity, heading structure, list structure,
table structure, code block score, paragraph structure, and word
overlap. Processing time was also recorded per page.
Results are reported as mean scores per strategy across the full
90-page corpus.

\section{Overall Strategy Performance (RQ1, RQ2)}

Table~\ref{tab:results_overview} shows the mean scores for all four
strategies across the full corpus.

\begin{table}[ht]
\centering
\caption{Mean metric scores per strategy across 90 pages (3 documents)}
\label{tab:results_overview}
\renewcommand{\arraystretch}{1.4}
\small\setlength{\tabcolsep}{4pt}
\begin{tabularx}{\textwidth}{lXXXXXXX}
\toprule
\textbf{Strategy} & \textbf{Text sim.} & \textbf{Heading} & \textbf{Lists} & \textbf{Tables} & \textbf{Code} & \textbf{Paras.} & \textbf{W. overlap} \\
\midrule
Text-only & \textbf{0.860} & \textbf{0.565} & 0.859 & 0.922 & \textbf{0.889} & \textbf{0.726} & 0.761 \\
Adaptive  & 0.838 & 0.542 & 0.868 & \textbf{0.956} & 0.752 & 0.636 & 0.794 \\
Image-only & 0.827 & 0.429 & \textbf{0.914} & 0.944 & 0.250 & 0.156 & \textbf{0.811} \\
Hybrid    & 0.818 & 0.348 & 0.891 & 0.944 & 0.736 & 0.597 & 0.795 \\
\bottomrule
\end{tabularx}
\end{table}

\begin{figure}[ht]
    \centering
    \includegraphics[width=\textwidth]{performance.png}
    \caption{Mean metric scores per strategy across 90 pages}
    \label{fig:performance}
\end{figure}

The most striking result is that text-only scores highest on five of the
seven metrics, including text similarity (0.860) and code blocks (0.889).
This is a non-obvious finding. The strategy uses no LLM and costs under
200\,ms per page. Its strength comes from the document corpus: thesis
documents consist primarily of continuous prose, and PyMuPDF's text
extraction is very accurate on such content.

The adaptive strategy scores highest on table detection (0.956), which is
0.034 points above text-only and 0.012 above image-only. This reflects its
design: pages classified as TABLE or MIXED are routed to the vision model,
which handles table layouts better than direct text extraction.
On the remaining metrics, adaptive is consistently second or close to
second among the four strategies.

Image-only shows a severe collapse on paragraph structure (0.156) and code
blocks (0.250). Given only a rendered page image and no extracted text, the model does not
reproduce the paragraph or code block structure from the reference. The model paraphrases and reflows the
content rather than preserving it verbatim. This makes image-only unsuitable
for documents where structural fidelity matters.

Hybrid combines image input with extracted text, which partially recovers
the paragraph and code block scores (0.597 and 0.736 respectively) compared
to image-only. However, it scores lowest on heading structure (0.348) and
text similarity (0.818). Providing both modalities does not produce the best
result on any single metric.

These results partially answer RQ1 and RQ2. The adaptive strategy does not
outperform text-only on most metrics, which contradicts hypothesis H1 for
text-heavy documents. It does outperform all strategies on table structure.
On efficiency and structural robustness across all page types, the adaptive
design shows a clear advantage, which Section~\ref{sec:efficiency} details.

\section{Efficiency (RQ3)}
\label{sec:efficiency}

Table~\ref{tab:efficiency} shows the mean processing time per page for
each strategy.

\begin{table}[ht]
\centering
\caption{Mean processing time per page by strategy}
\label{tab:efficiency}
\renewcommand{\arraystretch}{1.4}
\begin{tabular}{lrr}
\toprule
\textbf{Strategy} & \textbf{Avg time (ms)} & \textbf{Relative to text-only} \\
\midrule
Text-only  &    197.3 & $1\times$ \\
Adaptive   & 20{,}790.1 & $105\times$ \\
Image-only & 64{,}696.3 & $328\times$ \\
Hybrid     & 68{,}009.5 & $345\times$ \\
\bottomrule
\end{tabular}
\end{table}

\begin{figure}[ht]
    \centering
    \includegraphics[width=0.85\textwidth]{efficiency.png}
    \caption{Mean processing time per page by strategy (log scale)}
    \label{fig:efficiency}
\end{figure}

Text-only is by far the fastest strategy at 197\,ms per page on average.
It makes no LLM calls, so processing time is bounded by disk read and
PyMuPDF parsing speed. The adaptive strategy averages 20{,}790\,ms per page.
This is 105 times slower than text-only but 3.1 times faster than
image-only and 3.3 times faster than hybrid. The speedup comes directly
from the page classifier: TEXT pages skip the LLM entirely, and only pages
classified as IMAGE, TABLE, FORMULA, or MIXED go through the vision model.

Image-only and hybrid are the slowest strategies at 64{,}696\,ms and
68{,}009\,ms respectively. Every page is rendered as a base64-encoded image
and sent to the vision-language model, regardless of whether the page
contains any visual content. Hybrid is slightly slower than image-only
because it additionally sends the extracted text alongside the image,
increasing the input payload on every request. This extra cost does not
translate into a consistent accuracy improvement.

Token usage tells the same story from a different angle.
Table~\ref{tab:tokens} shows the total tokens consumed across all 90 pages
and the average per LLM call.

\begin{table}[ht]
\centering
\caption{Token usage per strategy across 90 pages}
\label{tab:tokens}
\renewcommand{\arraystretch}{1.4}
\begin{tabular}{lrrrr}
\toprule
\textbf{Strategy} & \textbf{LLM calls} & \textbf{Avg tokens / call} & \textbf{Total tokens} \\
\midrule
Text-only  & 0  & ---     & 0       \\
Adaptive   & 23 & 3{,}986 & 91{,}678  \\
Image-only & 90 & 3{,}516 & 316{,}440 \\
Hybrid     & 90 & 3{,}970 & 357{,}300 \\
\bottomrule
\end{tabular}
\end{table}

The per-call token cost is similar for all three VLM strategies, ranging
from 3{,}516 to 3{,}986 tokens per call. The difference is in how many
calls are made. Image-only and hybrid call the LLM on every one of the
90 pages. Adaptive calls it on only 23 pages --- the pages the classifier
identified as IMAGE, FORMULA, TABLE, or MIXED. The other 67 pages are
TEXT pages that go through pymupdf4llm with no LLM call at all.
This reduces total token consumption by 71\% compared to image-only and
74\% compared to hybrid. Hybrid consumes the most tokens of any strategy
because it sends both the rendered page image and the extracted text in
every request, which increases the input length on every call.

These results confirm hypothesis H3: the adaptive system is significantly
cheaper than image-only or hybrid in both processing time and token usage,
while maintaining comparable or better accuracy. The efficiency gain grows
with document length --- a 300-page document would consume roughly
$300 \times 3{,}986 \approx 1.2$\,million tokens with image-only but only
around $80{,}000$ with adaptive, assuming a similar content distribution.

\section{Per-Page Analysis and Consistency (RQ5)}

The per-page breakdown shows that no single strategy wins every page.
Across the 30 evaluated page positions, hybrid produces the best text
similarity on 14 pages, adaptive on 11, image-only on 5, and text-only
on 2. The variation reflects the content mix within each document.
Pages with complex table or figure layouts favour hybrid and adaptive.
Dense text pages favour text-only.

Page 6 is the lowest-scoring page across all strategies, with a best
similarity of 0.728. Inspection of the reference file shows that this
page contains a table of contents with dot-leader formatting, which none
of the strategies reproduce cleanly. This is a known limitation of both
text extraction (dot leaders appear as fragmented text) and vision
extraction (the model does not reproduce the exact list format).

The strategy rankings are consistent across the three thesis documents.
Text-only leads on text similarity in all three. Adaptive leads on table
structure in all three. This consistency addresses RQ5: the relative
performance ordering does not change across documents, even though absolute
scores vary slightly depending on content density. A document with more
formula or image pages would widen the gap between adaptive and text-only
because the classifier would route more pages to the vision model.

\section{Output Quality Metrics (RQ4)}

The seven-metric framework reveals failure modes that a single similarity
score would hide. Image-only scores 0.827 on text similarity, which looks
acceptable. But its paragraph structure score of 0.156 and code block score
of 0.250 expose a fundamental problem: the model reproduces most of the
words but loses the structural layout entirely. A RAG system ingesting this
output would receive the right vocabulary in the wrong format, which breaks
chunk boundaries and damages retrieval precision.

The heading structure scores for all four strategies are below 0.600. This
is a consistent weakness. The strategies differ in how they assign heading
levels: text-only derives headings from font size heuristics, while the
vision strategies rely on the model to infer structure from the visual
layout. Neither approach perfectly matches the reference heading hierarchy.
This suggests that heading level detection is the most difficult structural
task across all evaluated strategies.

Table structure is the metric where the strategies diverge most clearly.
Adaptive scores 0.956, while text-only scores 0.922 and both image-only
and hybrid score 0.944. The 0.034 gap between adaptive and text-only is
meaningful: on pages where text extraction produces malformed or missing
table markup, the vision model restores it. This is the metric where
routing to a VLM provides a measurable benefit.

Word overlap scores are close across all four strategies, ranging from
0.761 (text-only) to 0.811 (image-only). This means all strategies
recover most of the correct vocabulary. The differences between strategies
are in structure and formatting, not in which words appear in the output.
